% Datasets
\section{Dataset Preparation and Harmonization}
\label{sec:datasets}

This study uses the KITTI Object Detection dataset~\cite{geiger2012kitti} as the primary
source for model training and in-domain validation, while a representative subset of the
Waymo Open Dataset~\cite{sun2020waymo} is used exclusively for external cross-dataset
evaluation. This design allows the evaluated detectors to be compared under both familiar
and shifted driving conditions without using Waymo data for training, fine-tuning,
hyperparameter selection, or model selection. Table~\ref{tab:dataset_comparison}
summarizes the composition and roles of the three frozen partitions, and
Table~\ref{tab:class_mapping} records the class-harmonization rules.

\begin{table}[t]
    \centering
    \caption{Dataset composition and experimental roles of the three frozen partitions.}
    \label{tab:dataset_comparison}
    \begin{tabular}{lcccc}
        \toprule
        Metric & KITTI (all) & KITTI train & KITTI val. & Waymo ext. \\
        \midrule
        Experimental role & Source labeled benchmark & Model training & In-domain validation & External validation only \\
        Source split & Official labeled training set & Project train split & Project validation split & Official Waymo validation split \\
        Images & 7481 & 5985 & 1496 & 996 \\
        Driving segments & --- & --- & --- & 25 \\
        Target boxes & 39086 & 31294 & 7792 & 24819 \\
        Vehicle boxes & 32750 & 26278 & 6472 & 16928 \\
        Pedestrian boxes & 4709 & 3729 & 980 & 7127 \\
        Cyclist boxes & 1627 & 1287 & 340 & 764 \\
        Images with no target boxes & 0 & 0 & 0 & 12 \\
        Camera & Left color camera (image\_2) & Left color camera (image\_2) & Left color camera (image\_2) & FRONT \\
        Sampling rule & Use all official labeled images & Frozen stratified split & Frozen stratified split & every\_5th\_front\_frame\_starting\_at\_first \\
        Random seed & --- & 42 & 42 & --- \\
        \bottomrule
    \end{tabular}
\end{table}

\begin{table}[t]
    \centering
    \caption{Class harmonization mapping KITTI and Waymo annotations into the common Vehicle, Pedestrian, and Cyclist task.}
    \label{tab:class_mapping}
    \begin{tabular}{llll}
        \toprule
        Dataset & Source class & Action & Unified class \\
        \midrule
        KITTI & Car & map & Vehicle \\
        KITTI & Van & map & Vehicle \\
        KITTI & Truck & map & Vehicle \\
        KITTI & Pedestrian & map & Pedestrian \\
        KITTI & Person\_sitting & map & Pedestrian \\
        KITTI & Cyclist & map & Cyclist \\
        KITTI & Tram & ignore & --- \\
        KITTI & Misc & ignore & --- \\
        KITTI & DontCare & ignore & --- \\
        Waymo & Vehicle & map & Vehicle \\
        Waymo & Pedestrian & map & Pedestrian \\
        Waymo & Sign & ignore & --- \\
        Waymo & Cyclist & map & Cyclist \\
        \bottomrule
    \end{tabular}
\end{table}


\subsection{KITTI Dataset Preparation}
\label{sec:datasets_kitti}

The KITTI object-detection benchmark was selected as the primary labeled dataset. The
official object-detection training set contains $7{,}481$ left-color camera images stored
in the \texttt{image\_2} directory, together with $7{,}481$ corresponding annotation files
in \texttt{label\_2} and $7{,}481$ camera-calibration files. In KITTI terminology,
\texttt{image\_2} denotes the left color camera stream, and the corresponding camera
projection matrix is represented by $P_2$ in the calibration files.

The official KITTI testing set was not included in the local experimental evaluation
because its ground-truth annotations are not publicly available. Consequently, the
official labeled training set was divided into project-specific training and validation
partitions. Before splitting, all images, annotation files, and calibration files were
subjected to an integrity-verification procedure, which confirmed that all $7{,}481$
samples had matching image, label, and calibration identifiers; that all images were
readable; that all $51{,}865$ annotation rows had valid structures; and that no invalid
bounding boxes were found.

Special handling was applied to KITTI \texttt{DontCare} annotations, which may contain
placeholder truncation and occlusion values of $-1$ that are valid for ignored regions and
were therefore not treated as annotation errors. \texttt{DontCare} regions were preserved
for later evaluation handling but were not treated as a target object class. The verified
KITTI dataset contains $51{,}865$ original annotation rows, with an original class
distribution of $28{,}742$ Car, $2{,}914$ Van, $1{,}094$ Truck, $4{,}487$ Pedestrian,
$222$ Person\_sitting, $1{,}627$ Cyclist, $511$ Tram, $973$ Misc, and $11{,}295$
DontCare annotations.

A fixed stratified split was created using a random seed of $42$, preserving the presence
of the Vehicle, Pedestrian, and Cyclist classes in addition to the distribution of mapped
target-object density. This produced $5{,}985$ training images and $1{,}496$ in-domain
validation images. The training partition contains $26{,}278$ Vehicle boxes, $3{,}729$
Pedestrian boxes, and $1{,}287$ Cyclist boxes; the validation partition contains
$6{,}472$ Vehicle boxes, $980$ Pedestrian boxes, and $340$ Cyclist boxes. The same frozen
split is used for all detectors to ensure a fair comparison.

\subsection{Waymo Representative External-Validation Subset}
\label{sec:datasets_waymo}

The Waymo Open Dataset was used to construct an external-validation subset representing a
broader range of locations, traffic densities, illumination conditions, and weather
conditions, considering only the official Waymo validation split. The selection process
began with metadata from all $202$ validation segments, analyzing time of day, location,
weather, object counts, and traffic-density characteristics for each segment. The original
catalog contained $160$ daytime, $23$ dawn/dusk, and $19$ nighttime segments; $93$
Phoenix, $88$ San Francisco, and $21$ other-location segments; and $201$ sunny versus a
single rainy segment.

A broad candidate pool of $70$ segments was first selected to maximize coverage of rare
and important conditions, including all available nighttime and dawn/dusk segments, the
only rainy segment, pedestrian-rich scenes, cyclist-rich scenes, and daytime scenes
balanced across location and density groups. The 2D camera annotations of the candidates
were then analyzed using only the Waymo FRONT camera, and an optimization-based selection
procedure produced a final frozen set of $25$ representative driving segments. The final
segment set contains $14$ daytime, six dawn/dusk, and five nighttime segments; $12$ San
Francisco, $11$ Phoenix, and two other-location segments; and ten high-density, eight
low-density, and seven medium-density segments, with pedestrians in $22$ segments and
cyclists in $18$ segments.

After the segment list was frozen, the corresponding camera images and calibration files
were downloaded. To reduce temporal redundancy, the FRONT-camera frames in each segment
were sorted chronologically and sampled uniformly by retaining every fifth frame, yielding
$996$ images. The selected images were matched with their 2D camera boxes using the
segment identifier, frame timestamp, and camera identifier; only Vehicle, Pedestrian, and
Cyclist annotations were retained, with Sign annotations excluded. The final Waymo subset
contains $24{,}819$ target boxes ($16{,}928$ Vehicle, $7{,}127$ Pedestrian, and $764$
Cyclist), and $12$ images with no target boxes are retained intentionally as negative
samples for false-positive analysis.

The Waymo subset is used only for external evaluation and is not used for model training,
retraining, fine-tuning, validation-based checkpoint selection, threshold selection, or
hyperparameter optimization. This separation prevents information leakage and allows the
external results to reflect direct cross-dataset generalization.

\subsection{Class Harmonization}
\label{sec:datasets_harmonization}

To enable a consistent comparison, both datasets were mapped into a common three-class
detection task---Vehicle, Pedestrian, and Cyclist. For KITTI, Car, Van, and Truck were
merged into Vehicle; Pedestrian and Person\_sitting were merged into Pedestrian; and
Cyclist was mapped directly to Cyclist. Tram, Misc, and DontCare were excluded from the
target classes, with DontCare retained as an ignored region where appropriate. For Waymo,
Vehicle, Pedestrian, and Cyclist were mapped directly to their corresponding unified
classes, while Sign annotations were ignored (Table~\ref{tab:class_mapping}).

Following harmonization, KITTI contains $39{,}086$ target boxes ($32{,}750$ Vehicle,
$4{,}709$ Pedestrian, and $1{,}627$ Cyclist), with a further $12{,}779$ annotations
retained as ignored. The Waymo representative subset contains $24{,}819$ target boxes
under the same three-class definition. This harmonization ensures that all detector
families are trained and assessed using an identical class order and semantic task
definition.

\subsection{Dataset Statistics and Visual Verification}
\label{sec:datasets_verification}

Image-level and object-level statistics were generated for the complete KITTI dataset and
for its training and validation partitions, including original and harmonized class
distributions, object density, bounding-box dimensions, normalized area, occlusion,
truncation, and descriptive KITTI difficulty categories. Bounding boxes were additionally
divided into small, medium, and large descriptive groups according to the lower and upper
thirds of the normalized target-box area distribution; these project-specific groups are
used for dataset analysis and do not replace the official KITTI difficulty definitions.
Visual-inspection samples were generated for Vehicle, Pedestrian, Cyclist, mixed, crowded,
occluded, truncated, and DontCare scenes, with equivalent checks for the Waymo subset,
confirming that the generated annotations corresponded to the source labels and that class
harmonization was applied correctly. Overall, the final protocol consists of $5{,}985$
KITTI images for training, $1{,}496$ KITTI images for in-domain validation, and $996$
Waymo FRONT-camera images for external cross-dataset evaluation.
